\subsection*{(3.3) The general case}

\begin{theorem}[3.3.1]
Let $f:X\to Y$ be a morphism of schemes of finite type over $\mathbb Z$, and let $\mathcal F$ be a sheaf on $X$. If $\mathcal F$ is mixed of weights $\le n$, then, for each $i$, $R^if_{!}\mathcal F$ is mixed of weights $\le n+i$.
\end{theorem}

The proof uses successive devissages. Exact sequences reduce the assertion in the coefficient sheaf; an open-closed decomposition reduces it in $X$; base change to an open part and its complement reduces it in $Y$; Leray spectral sequences give transitivity. Universal homeomorphisms do not affect etale cohomology. Finally, morphisms of relative dimension zero preserve pointwise purity and have no higher direct images. These reductions reduce to the case where $f$ has relative dimension one, $\mathcal F$ is lisse and pointwise pure, and $X$ is the complement in a smooth projective relative curve of a divisor finite etale over $Y$, with tame ramification along the divisor.

On each geometric fiber, Theorem (3.2.3) gives purity for $R^if_{!}j_{*}\mathcal F$. The local-monodromy theorem (1.8.4) gives the weights of the graded pieces of the local monodromy filtration on $i^*j_{*}\mathcal F$. The finite part is handled by the relative-dimension-zero case, and the exact sequence
\[
0\to j_{!}\mathcal F\to j_{*}\mathcal F\to i_{*}i^*j_{*}\mathcal F\to0
\]
finishes the proof.

\paragraph{(3.3.2).}
Lower bounds for weights are obtained from the $p$-adic estimates of SGA 7 XXI, 5. A sheaf is called integral if all Frobenius eigenvalues are algebraic integers. The weights of a mixed integral sheaf are necessarily non-negative, since the norm of a pure algebraic integer of weight $n$ has absolute value $q^{dn/2}$.

\begin{corollary}[3.3.3]
Let $f:X\to Y$ be of finite type over $\mathbb Z$, and let $\mathcal F$ be an integral sheaf on $X$. If $\mathcal F$ is mixed of weights $\le n$, then $R^if_{!}\mathcal F$ is mixed of weights between $0$ and $n+i$. If $i$ is larger than the maximum dimension $d$ of the fibers, the weights lie between $2(i-d)$ and $n+i$.
\end{corollary}

For $Y=\operatorname{Spec}\Fq$ this becomes:

\begin{corollary}[3.3.4]
Let $X_0$ be of finite type over $\Fq$, and let $\mathcal F_0$ be mixed of weights $\le n$ on $X_0$. Then $H^i_c(X,\mathcal F)$ is mixed of weights $\le n+i$. If $\mathcal F_0$ is integral, the weights are non-negative; if moreover $i>d=\dim X$, they are $\ge2(i-d)$.
\end{corollary}

\begin{corollary}[3.3.5]
If $X_0$ is smooth of finite type over $\Fq$ and $\mathcal F_0$ is lisse mixed of weights $\ge n$, then $H^i(X,\mathcal F)$ is mixed of weights $\ge n+i$.
\end{corollary}
This is Poincare duality applied to the dual sheaf. Hence:

\begin{corollary}[3.3.6]
If $X_0$ is smooth of finite type over $\Fq$ and $\mathcal F_0$ is lisse pure of weight $n$, then the image of
\[
H^i_c(X,\mathcal F)\longrightarrow H^i(X,\mathcal F)
\]
is pure of weight $n+i$.
\end{corollary}

\paragraph{(3.3.7)--(3.3.9).}
If an algebraic number $\alpha$ is pure of weight $n$ relative to $q$, normalize $p$-adic valuations by $v(q)=1$ and attach to $\alpha$ the pair $(v(\alpha),v(q^n\alpha^{-1}))$. For integral eigenvalues both coordinates are non-negative and sum to $n$. For the constant sheaf, the preceding bounds localize these pairs in the corresponding triangles in the $(r,s)$-plane. In particular, for $X_0$ proper and smooth one obtains again the main theorem of I, with projective replaced by proper:
\[
\det(1-Ft,H^i(X,\Ql))\in\mathbb Z[t]
\]
is independent of $\ell\ne p$, and all complex roots have absolute value $q^{i/2}$.

\paragraph{(3.3.10)--(3.3.11).}
The theorem remains valid with ``mixed'' replaced by ``$\iota$-mixed'' and real-valued weights. The only weights in $R^if_{!}\mathcal F$ are those $\le\beta+i$ congruent modulo $\mathbb Z$ to a weight already occurring in $\mathcal F$. The smoothness assumption is used only to invoke Poincare duality, and may therefore be replaced by the corresponding rational-homology-manifold condition
\[
Ra^!\Ql=\Ql(n)[2n]
\]
for the structural morphism $a:X_0\to\Spec\Fq$.

